% ======================================================================
%  preamble.tex — оформление пособия «Квантовая физика: ответы на билеты»
%  Компиляция: XeLaTeX (см. README.md)
% ======================================================================

% ---------- Геометрия страницы ----------
\usepackage[a4paper,top=22mm,bottom=24mm,inner=24mm,outer=22mm,
            headheight=14pt,headsep=8mm,footskip=11mm]{geometry}

% ---------- Математика (до unicode-math) ----------
\usepackage{amsmath}
\usepackage{amssymb}
\usepackage{mathtools}

% ---------- Шрифты и язык ----------
\usepackage{fontspec}
\usepackage{unicode-math}
\usepackage{polyglossia}
\setdefaultlanguage{russian}
\setotherlanguage{english}

% Computer Modern Unicode — полноценная кириллица, гармонирует с Latin Modern Math.
% Файлы шрифтов лежат в каталоге fonts/ проекта (компилировать из корня проекта).
\setmainfont{cmu-serif}[
  Path           = fonts/,
  Extension      = .ttf,
  UprightFont    = *-500-roman,
  ItalicFont     = *-500-italic,
  BoldFont       = *-700-roman,
  BoldItalicFont = *-700-italic,
  Ligatures      = TeX]
\setsansfont{cmu-sans-serif}[
  Path           = fonts/,
  Extension      = .ttf,
  UprightFont    = *-500-roman,
  ItalicFont     = *-500-italic,
  BoldFont       = *-700-roman,
  BoldItalicFont = *-700-italic,
  Ligatures      = TeX]
\setmonofont{cmu-typewriter-text-500-roman}[Path=fonts/, Extension=.ttf]
\newfontfamily\cyrillicfont{cmu-serif}[
  Path=fonts/, Extension=.ttf, UprightFont=*-500-roman, ItalicFont=*-500-italic,
  BoldFont=*-700-roman, BoldItalicFont=*-700-italic, Ligatures=TeX]
\newfontfamily\cyrillicfontsf{cmu-sans-serif}[
  Path=fonts/, Extension=.ttf, UprightFont=*-500-roman, ItalicFont=*-500-italic,
  BoldFont=*-700-roman, BoldItalicFont=*-700-italic, Ligatures=TeX]
\newfontfamily\cyrillicfonttt{cmu-typewriter-text-500-roman}[Path=fonts/, Extension=.ttf]

\setmathfont{latinmodern-math.otf}

\usepackage{icomma}           % десятичная запятая в формулах: 0,529
\linespread{1.06}
\setlength{\parindent}{1.25em}
\setlength{\parskip}{0.25em plus 0.15em}
\setlength{\emergencystretch}{2.5em}
\clubpenalty=10000
\widowpenalty=10000
\displaywidowpenalty=10000

% ---------- Цвета ----------
\usepackage{xcolor}
\definecolor{accent}{HTML}{1F4E79}
\definecolor{accentlight}{HTML}{EAF1F8}
\definecolor{defcol}{HTML}{2E6DA4}
\definecolor{defbg}{HTML}{F2F7FB}
\definecolor{keycol}{HTML}{8A3B12}
\definecolor{keybg}{HTML}{FBF4EE}
\definecolor{dercol}{HTML}{4F6B2E}
\definecolor{derbg}{HTML}{F5F8F0}
\definecolor{remcol}{HTML}{6B6B6B}
\definecolor{examcol}{HTML}{6A1B5A}
\definecolor{exambg}{HTML}{F8F0F6}
\definecolor{qbg}{HTML}{F1F1F1}

% ---------- Графика ----------
\usepackage{graphicx}
\usepackage{tikz}
\usetikzlibrary{arrows.meta,calc,decorations.pathmorphing,decorations.markings,
                patterns,positioning,shapes.geometric,angles,quotes}
\usepackage{pgfplots}
\pgfplotsset{compat=1.18,
  every axis/.append style={font=\small, line width=0.6pt,
     tick style={line width=0.5pt}},
  every axis label/.append style={font=\small}}
\usepgfplotslibrary{fillbetween}
\tikzset{>={Stealth[length=2.2mm]}}

% ---------- Таблицы, списки, подписи ----------
\usepackage{booktabs}
\usepackage{array}
\usepackage{multirow}
\usepackage{tabularx}
\usepackage{enumitem}
\setlist{itemsep=1pt, topsep=3pt, parsep=1pt, leftmargin=1.6em}
\usepackage[font=small,labelfont=bf,labelsep=period,
            figurename=Рис.,tablename=Таблица,skip=4pt]{caption}
\usepackage{float}
\usepackage{needspace}

% ---------- Заголовки ----------
\usepackage{titlesec}
\titleformat{\chapter}[display]
  {\normalfont\sffamily\color{accent}}
  {\large\MakeUppercase{Билет}~\thechapter}
  {4pt}
  {\LARGE\bfseries\raggedright}
  [\vspace{3pt}{\color{accent}\titlerule[0.8pt]}]
\titlespacing*{\chapter}{0pt}{-6pt}{10pt}
\titleformat{name=\chapter,numberless}[display]
  {\normalfont\sffamily\color{accent}}{}{0pt}
  {\LARGE\bfseries\raggedright}[\vspace{3pt}{\color{accent}\titlerule[0.8pt]}]
\titlespacing*{name=\chapter,numberless}{0pt}{-6pt}{10pt}

\titleformat{\section}
  {\normalfont\sffamily\large\bfseries\color{accent}}{\thesection}{0.7em}{}
\titlespacing*{\section}{0pt}{2.0ex plus 0.6ex minus 0.3ex}{1.0ex plus 0.3ex}
\titleformat{\subsection}
  {\normalfont\sffamily\bfseries}{\thesubsection}{0.6em}{}
\titlespacing*{\subsection}{0pt}{1.6ex plus 0.5ex minus 0.3ex}{0.6ex plus 0.2ex}
\titleformat{\subsubsection}[runin]
  {\normalfont\bfseries}{}{0pt}{}[.]
\titlespacing*{\subsubsection}{0pt}{1.0ex plus 0.3ex}{0.6em}

% Запрет «висячих» заголовков внизу страницы
\usepackage{etoolbox}
\pretocmd{\section}{\needspace{12\baselineskip}}{}{}
\pretocmd{\subsection}{\needspace{10\baselineskip}}{}{}

% ---------- Оглавление ----------
\usepackage{titletoc}
\titlecontents{chapter}[5.4em]
  {\addvspace{3pt}\sffamily}
  {\contentslabel[\textbf{Билет~\thecontentslabel}]{5.4em}}
  {}
  {\titlerule*[0.6pc]{.}\contentspage}
\setcounter{tocdepth}{0}

% ---------- Колонтитулы ----------
\usepackage{fancyhdr}
\pagestyle{fancy}
\fancyhf{}
\renewcommand{\chaptermark}[1]{\markboth{Билет~\thechapter.\ #1}{}}
\fancyhead[L]{\small\sffamily\color{accent}\leftmark}
\fancyhead[R]{}
\fancyfoot[C]{\small\sffamily\thepage}
\renewcommand{\headrulewidth}{0.4pt}
\fancypagestyle{plain}{\fancyhf{}\fancyfoot[C]{\small\sffamily\thepage}
  \renewcommand{\headrulewidth}{0pt}}

% ---------- Цветные блоки ----------
\usepackage[most]{tcolorbox}
\tcbset{
  basebox/.style={enhanced, breakable, boxrule=0pt, arc=1.5pt, outer arc=1.5pt,
     left=7pt, right=7pt, top=4pt, bottom=4pt, before skip=7pt, after skip=7pt,
     fonttitle=\sffamily\bfseries\small, coltitle=black, attach title to upper,
     after title={.\ }, parbox=false}
}
% Формулировка билета
\newtcolorbox{qbox}{enhanced, breakable, colback=qbg, colframe=accent,
  boxrule=0pt, leftrule=3pt, arc=0pt, left=8pt, right=8pt, top=6pt, bottom=6pt,
  before skip=4pt, after skip=10pt, fontupper=\itshape, parbox=false}
% Определение
\newtcolorbox{definition}[1][]{basebox, colback=defbg, colframe=defcol,
  leftrule=2.5pt, title={Определение\ifstrempty{#1}{}{ (#1)}}}
% Закон / постулат / принцип / теорема
\newtcolorbox{principle}[1]{basebox, colback=defbg, colframe=defcol,
  leftrule=2.5pt, title={#1}}
% Ключевой результат
\newtcolorbox{keyresult}[1][Ключевой результат]{basebox, colback=keybg,
  colframe=keycol, boxrule=0.6pt, leftrule=2.5pt, title={#1}}
% Вывод (нумерация внутри билета; метка: \begin{derivation}[label=der:x]{Заголовок})
\newtcolorbox[auto counter, number within=chapter]{derivbox}[2][]{basebox,
  colback=derbg, colframe=dercol, leftrule=2.5pt,
  title={Вывод~\thetcbcounter: #2}, after title={\par\smallskip}, #1}
\newenvironment{derivation}[2][]
  {\par\needspace{9\baselineskip}\begin{derivbox}[#1]{#2}}
  {\end{derivbox}}
% Защита от отрыва заголовка блока от его содержимого
\BeforeBeginEnvironment{definition}{\par\needspace{4\baselineskip}}
\BeforeBeginEnvironment{principle}{\par\needspace{4\baselineskip}}
\BeforeBeginEnvironment{keyresult}{\par\needspace{4\baselineskip}}
\BeforeBeginEnvironment{remark}{\par\needspace{4\baselineskip}}
% Замечание
\newtcolorbox{remark}[1][Замечание]{basebox, colback=white, colframe=remcol,
  leftrule=1.2pt, title={#1}, fontupper=\small}
% Экзаменационный чек-лист
\newtcolorbox{examcheck}{enhanced, breakable, colback=exambg, colframe=examcol,
  boxrule=0.6pt, arc=1.5pt, left=8pt, right=8pt, top=6pt, bottom=6pt,
  before skip=6pt, after skip=6pt, parbox=false}

% ---------- Итоговые формулы ----------
% \keyformula{формула}{обозначения}{условия применимости}{что позволяет найти}
\newcommand{\keyformula}[4]{%
  \par\needspace{6\baselineskip}%
  \begin{tcolorbox}[enhanced, breakable=false, colback=keybg, colframe=keycol,
     boxrule=0.5pt, leftrule=2.5pt, arc=1.5pt, left=7pt, right=7pt,
     top=4pt, bottom=4pt, before skip=5pt, after skip=5pt,
     before upper={\setlength{\parindent}{0pt}\setlength{\parskip}{1pt}}]
  {\centering $\displaystyle #1$\par}\vspace{5pt}
  {\small\textbf{Обозначения:} #2\par
   \textbf{Условия:} #3\par
   \textbf{Смысл:} #4\par}
  \end{tcolorbox}}

% Подзаголовки шаблона внутри «Что необходимо знать на экзамене»
\newcommand{\checkhead}[1]{\par\smallskip\noindent{\sffamily\bfseries\color{examcol}#1}\par\nopagebreak}

% ---------- Макросы обозначений ----------
\newcommand{\vect}[1]{\symbf{#1}}            % вектор
\newcommand{\op}[1]{\hat{#1}}                 % оператор
\newcommand{\vop}[1]{\hat{\symbf{#1}}}        % векторный оператор
\newcommand{\dd}{\mathrm{d}}                  % дифференциал
\newcommand{\ii}{\mathrm{i}}                  % мнимая единица
\newcommand{\ee}{\mathrm{e}}                  % основание экспоненты
\newcommand{\diff}[2]{\frac{\dd #1}{\dd #2}}
\newcommand{\pdiff}[2]{\frac{\partial #1}{\partial #2}}
\DeclarePairedDelimiter{\abs}{\lvert}{\rvert}
\DeclarePairedDelimiter{\norm}{\lVert}{\rVert}
\DeclarePairedDelimiter{\avg}{\langle}{\rangle}
\DeclarePairedDelimiterX{\comm}[2]{[}{]}{#1,#2}
\newcommand{\const}{\mathrm{const}}
\newcommand{\sh}{\operatorname{sh}}
\newcommand{\ch}{\operatorname{ch}}
\renewcommand{\th}{\operatorname{th}}
\newcommand{\tg}{\operatorname{tg}}
\newcommand{\ctg}{\operatorname{ctg}}
\newcommand{\arctg}{\operatorname{arctg}}
\newcommand{\Ry}{\mathrm{Ry}}
\newcommand{\lC}{\lambda_{\mathrm C}}
\newcommand{\aB}{a_{\mathrm B}}
\newcommand{\muB}{\mu_{\mathrm B}}
\newcommand{\mel}{m_e}
\newcommand{\unit}[1]{\,\mathrm{#1}}
\newcommand{\eV}{\,\text{эВ}}
\newcommand{\keV}{\,\text{кэВ}}
\newcommand{\MeV}{\,\text{МэВ}}
\newcommand{\GeV}{\,\text{ГэВ}}
\newcommand{\fm}{\,\text{фм}}
\newcommand{\nm}{\,\text{нм}}
\newcommand{\pmet}{\,\text{пм}}
\newcommand{\Ang}{\,\text{Å}}
\newcommand{\cm}{\,\text{см}}
\newcommand{\m}{\,\text{м}}
\newcommand{\s}{\,\text{с}}
\newcommand{\K}{\,\text{К}}
\newcommand{\Tl}{\,\text{Тл}}
\newcommand{\nuc}[3]{{}^{#2}_{#3}\mathrm{#1}}  % \nuc{U}{238}{92}
\newcommand{\nucl}[2]{{}^{#2}\mathrm{#1}}      % \nucl{U}{238}
\renewcommand{\le}{\leqslant}
\renewcommand{\ge}{\geqslant}
\renewcommand{\leq}{\leqslant}
\renewcommand{\geq}{\geqslant}

% Нумерация формул внутри билета
\numberwithin{equation}{chapter}
\numberwithin{figure}{chapter}
\numberwithin{table}{chapter}

% ---------- Гиперссылки ----------
\usepackage[unicode=true, bookmarksnumbered=true, bookmarksopen=true,
            bookmarksopenlevel=0, pdfpagelabels]{hyperref}
\hypersetup{colorlinks=true, linkcolor=accent, urlcolor=accent, citecolor=accent,
  pdftitle={Квантовая физика: подробные ответы на экзаменационные вопросы},
  pdfsubject={Общая физика: квантовая физика, МФТИ, V семестр},
  pdfkeywords={квантовая физика, МФТИ, экзамен}}
\usepackage{bookmark}
\bookmarksetup{depth=2}

% Подпись «Билет» для \ref на главы
\newcommand{\tref}[1]{билет~\ref{#1}}
\newcommand{\Tref}[1]{Билет~\ref{#1}}
